% !PRINTED 58
\weberpage{58}

\begin{center}
\S.\,15.

\gesperrt{Indices für eine Potenz von 2 als Modul}.
\end{center}

Ein ähnlicher Satz muss nun für die Potenzen $2^\lambda$ der Primzahl 2 aufgestellt werden, die sich, wie schon vorhin bemerkt, anders verhält.

Ist $\lambda = 1$, so ist jede durch 2 nicht theilbare Zahl $n \equiv 1 \pmod{2}$. Ist $\lambda = 2$, so ist $-1$ als primitive Wurzel von 4 aufzufassen, denn jede ungerade Zahl $n$ genügt einer der Congruenzen

\[
n \equiv (-1)^\alpha \pmod{4},
\]

worin $\alpha = 0$ oder $= 1$ ist.

Aber schon für den Modul 8 existirt keine primitive Wurzel mehr, d.\,h. keine Zahl, durch deren Potenzen sich alle Zahlclassen ungerader Zahlen nach dem Modul 8 darstellen lassen. Denn ist $g$ irgend eine ungerade Zahl, so sind unter den Potenzen von $g$ höchstens die beiden $1, g$ nach dem Modul 8 verschieden, weil immer $g^2 \equiv 1 \pmod{8}$ ist. Nimmt man also $g$ nicht congruent 1 und nicht congruent $-1 \pmod{8}$, also $g = 3$ oder $= 5$, so ist jede ungerade Zahl einer der vier Zahlen

\[
(-1)^\alpha g^\beta \qquad \alpha = 0, 1, \quad \beta = 0, 1
\]

nach dem Modul 8 congruent.

Die Anzahl der Classen ungerader Zahlen nach dem Modul $2^\lambda$ ist $2^{\lambda-1}$. Nun ist aber für jede ungerade Zahl $g$, falls $\lambda > 2$ ist,

\weq{1}{g^{2^{\lambda-2}} \equiv 1 \pmod{2^\lambda}.}

Denn nehmen wir (1) als richtig an und setzen demgemäss

\[
g^{2^{\lambda-2}} = 1 + h\, 2^\lambda,
\]

und erheben ins Quadrat, so folgt:

\[
g^{2^{\lambda-1}} \equiv 1 \pmod{2^{\lambda+1}}.
\]

Ist also die Formel (1) für $\lambda$ richtig, so ist sie es auch für $\lambda + 1$, und da sie für $\lambda = 3$ gilt, so gilt sie allgemein. Daraus folgt, dass unter den Potenzen einer ungeraden Zahl $g$ höchstens $2^{\lambda-2}$ nach dem Modul $2^\lambda$ verschiedene vorkommen können.

Andererseits folgt aber leicht für $g = 5$, dass $5^{2^{\lambda-2}}$ die niedrigste Potenz ist, die nach dem Modul $2^\lambda$ mit der Einheit congruent wird. Denn ist $5^e$ die niedrigste Potenz, die der Ein-
%CONT>>
